\section{Problem Statement}

The processing of non-linear signals inherently involves dealing with complex dynamic systems. In many contexts, audio recordings may suffer from unwanted non-linear distortion. These distortions modify the signal in a way that depends not only on the direct time-steps but often on the echo of the signal. \cite{Keuninckx2017RealTime}.

\vspace{0.5cm}

The core research questions that will guide our comparative analysis throughout this study are formulated as follows:

\begin{itemize}
    \item \textbf{To what extent can the Reservoir Computing paradigm satisfy the constraints of real-time audio processing regarding algorithmic complexity, latency, and predictive accuracy for modeling strongly non-linear audio distortion ?}
    \item \textbf{Among the existing RC architectures, which specific model provides the optimal trade-off to perform this task at a standard 44.1 kHz sampling rate, maintaining a latency strictly below 15 ms on standard computing hardware ?}
\end{itemize}


\subsection{Research Topic}

Even if reservoir computing topic is relevant in many application (real time signal distortion, AI, timeseries forecasting, chaotic series, etc.) we will here focus on the field of real time audio distortion.


